\begin{musica}{Porto Na Memória}{Lino Lobão}{Tuna de Engenharia da Universidade do Porto}
\nota{Introdução: G G C C (pausa) Em C7 F Fm C G G7 C}
\nota{\phantom{Introdução: }C Em C7 F Fm C G G7 C}

\ac{C}{Num} amplo ca\ac{G}{sa}rio adorme\ac{C}{ci}do,
No pôr-do-sol que \ac{Cdim7}{bri}lha junto à \ac{Dm}{Foz},
O Porto canta a\ac{Am}{le}gre envaide\ac{G}{ci}do
A his\ac{Dm}{tó}ria que per\ac{G}{ten}ce a todos \ac{C}{nós}.

Re\ac{C}{lem}bro outros \ac{G}{tem}pos com sau\ac{C}{da}de
P´las ruas ruas e \ac{Cdim7}{rue}las tão ve\ac{Dm}{lhi}nhas,
P’la Sé, pela Ri\ac{G}{bei}ra ou na Trin\ac{C}{da}de
\ac{Am}{Can}ta-se o \ac{Dm}{fa}do, co\ac{G}{mem}-se as iscas quen\ac{C}{ti}nhas.

\nota{Refrão:}
\ac{C}{Vou} até ao \ac{Em}{cu}bo, vi\ac{C7}{ra}do para o \ac{F}{rio}
\ac{Fm}{}_ Olha a outra margem como es\ac{C}{tá} bonita agora
E \ac{G}{fi}co eterna\ac{G7}{men}te com o \ac{C}{Por}to na memória.
\ac{C}{Sal}tam os mi\ac{Em}{ú}dos da \ac{C7}{pon}te para o \ac{F}{rio},
\ac{Fm}{}_ Do Duque da Ribeira a\ac{C}{in}da reza a história
E o \ac{G}{Por}to tão ve\ac{G7}{lhi}nho vai fi\ac{C}{can}do na me\ac{A7}{mó}ria,
E o \ac{Dm}{Por}to tão ve\ac{G}{lhi}nho vai fi\ac{C}{can}do na memória.

Com euforia o povo sai p’ra rua,
É festa popular, é o São João;
Sardinhas, caldo verde com fartura,
Cantam-se marchas, lança-se o balão.

Tão nobre é este povo da cidade,
Que do trabalho tira o seu viver.
Vencido põe um fim à mocidade
Que é tradição trabalhar até morrer.

\nota{Refrão}

\posicoes{C: x32010, C7: x32310, Cdim7: x34242, Dm: xx0231, Em: 022000, F: 133211,
  Fm: 133111, G: 320033, Am: x02210}
\end{musica}
